\documentclass[12pt]{article}
\usepackage[utf8]{inputenc}
\usepackage[T1]{fontenc}
\usepackage{hyperref}
\title{Kenyah, Mainstream Language Resources for K-12 Curriculum Development}
\author{Compiled by Claude Sonnet 5.5 (claude-sonnet-5-5) for Andr\'as Kornai}
\date{}
\begin{document}
\maketitle

\section{Introduction}
Mainstream Kenyah (ISO 639-3 xkl; Glottocode main1275) is an Austronesian dialect cluster of North and East Kalimantan, Indonesia, and Sarawak, Malaysia. Wikipedia gives about 50,000 speakers. The search found overview material and audio listings, but no dictionary, grammar or school material.

\section{Primary Educational Resources}
None found.

\section{Linguistic and Reference Materials}
Wikipedia overview with dialect clusters (Lepo' Tau, Lepo' Ke, Lepo' Maut, Bakung) and phonology: \url{https://en.wikipedia.org/wiki/Mainstream_Kenyah_language}

\section{Religious and Community Resources}
No scripture translation or community organisation was identified in the search results. Global Recordings Network pages may link to gospel recordings, but this was not verified.

\section{Digital and Audio Resources}
Global Recordings Network lists the language and many named varieties:
\begin{itemize}
\item \url{https://globalrecordings.net/en/language/xkl}
\item Lepo' Tau: \url{https://www.globalrecordings.net/en/language/4137}
\item Lepo' Ke: \url{https://globalrecordings.net/en/language/20050}
\item Long Patah: \url{https://globalrecordings.net/en/language/574}
\item Badeng: \url{https://globalrecordings.net/en/language/11829}
\end{itemize}

\section{Curriculum Development Context}
The cluster is internally diverse and has no known standard orthography or school programme. Any curriculum would first have to choose a variety or varieties. Documentation of the individual varieties appears thin in what was found.
\end{document}